\documentclass[11pt]{article}

\usepackage[T1]{fontenc}
\usepackage[utf8]{inputenc}
\usepackage{lmodern}
\usepackage[margin=1in]{geometry}
\usepackage{amsmath,amssymb,amsthm}
\usepackage{microtype}
\usepackage[hidelinks,pdftitle={The 4 x 4 Torus Has Unbounded-Speed Robber Cop Number Four},pdfauthor={Xiaojian Zeng}]{hyperref}

\newtheorem{theorem}{Theorem}
\newtheorem{lemma}[theorem]{Lemma}
\newcommand{\Z}{\mathbb Z}
\newcommand{\supp}{\operatorname{supp}}

\title{The $4\times4$ Torus Has Unbounded-Speed Robber Cop Number Four}
\author{Xiaojian Zeng\\
  \small PBC School of Finance, Tsinghua University, Beijing 100190, China\\
  \small \texttt{xiaojianzeng@foxmail.com}}
\date{}

\begin{document}
\maketitle

\begin{abstract}
In the fast-robber game on a graph, the cops move simultaneously by at most
one edge and the robber then moves to any vertex reachable along a cop-free
path.  We determine the exact cop number of the smallest even square torus not
covered by the available asymptotic torus bounds:
\[
  c_{1,\infty}(C_4\mathbin{\square}C_4)=4.
\]
The proof is an exhaustive finite attractor certificate.  For four cops, a
rank-decreasing policy covers all 48,960 legal multiset states and forces
capture in at most four cop moves.  For three cops, a closed set of 4,576
states supplies an infinite evasion strategy.  A separately implemented
checker expands the multiset certificate to labelled cops and verifies every
ordered simultaneous move using exact integer and bit-mask operations.
\end{abstract}

\section{Context and result}

The fast-robber variant of Cops and Robbers was studied systematically for
bounded robber speed by Fomin, Golovach, Kratochv\'{\i}l, Nisse, and
Suchan~\cite{FominEtAl2010}.  We use the infinite-speed convention of
Kinnersley and Townsend~\cite{KinnersleyTownsend2022}: the cops choose their
initial vertices first, the robber then chooses an unoccupied vertex, and each
round consists of one simultaneous cop step followed, unless capture has
already occurred, by an arbitrarily long cop-free robber move.

For square tori, Kinnersley and Townsend proved
\[
  2m-24\leq c_{1,\infty}(C_m\mathbin{\square}C_m)\leq 2m
  \qquad(m\geq18)
\]
\cite[Theorem~2.7]{KinnersleyTownsend2022}.  Biswas's later
quotient-sweeping bound gives
\(c_{1,\infty}(C_m^{\square k})\leq2m^{k-1}\) for \(m\geq4\),
and restates the preceding two-dimensional estimate~\cite[Section~5.2]
{Biswas2026}.  Thus the cited results give the non-exact upper bound $8$ when
\(m=4\).  The contribution here is the exact value for this single
16-vertex graph.  It makes no assertion about larger tori and no priority
assertion beyond this bounded comparison with the cited literature.

Exact finite certificates have also been used to delimit small discrete
models in other settings, with the discovery computation separated from an
exact verifier that reconstructs the instance~\cite[Sections~6--7]
{ZengEtAl2026Certificates}.  We follow that reproducibility pattern here, but
the reachability-game reduction and all quantified checks needed for the
present theorem are stated below.

Let $G=C_4\mathbin{\square}C_4$, with vertex set $\Z_4^2$.  Distinct
vertices $(x,y)$ and $(x',y')$ are adjacent when exactly one coordinate is
equal and the other differs by $1$ modulo $4$.  Repeated cop positions,
cop collisions, swaps, and stays are allowed.  After the cops move to a set of
occupied vertices $S$, capture is tested immediately.  If the robber has not
been captured, its possible endpoints are exactly the vertices in the
component of $G-S$ containing its current vertex.  Vacated cop vertices may
therefore be used on that move.

\begin{theorem}\label{thm:main}
Under the preceding placement, move-order, and capture convention,
\[
  c_{1,\infty}(C_4\mathbin{\square}C_4)=4.
\]
\end{theorem}

\section{Finite attractor and certificates}

We first give the mathematical semantics of the computation.  A $k$-cop
configuration $C$ is a multiset of $k$ vertices.  Write $C\leadsto D$ if
the occurrences in $C$ can be matched bijectively to those in $D$ so that
each matched pair is equal or adjacent.  This is precisely a legal
simultaneous cop move, including moves between configurations with repeated
positions.  A legal state at the start of a cop turn is a pair $(C,r)$ with
$r\notin\supp(C)$.  If $r\notin\supp(D)$, let $R_D(r)$ be the vertex set
of the component of $r$ in $G-\supp(D)$.

Set $W_0=\varnothing$, and for $t\geq0$ define
\begin{align*}
W_{t+1}=W_t\ \cup\ 
\bigl\{(C,r):{}&\text{ there is a }D\text{ with }C\leadsto D\text{ such that}\\
&r\in\supp(D)\text{, or }(D,r')\in W_t
  \text{ for every }r'\in R_D(r)\bigr\}.
\end{align*}
This is the least reachability attractor: the cops choose $D$, after which
every legal robber response must lie in the earlier set.

\begin{lemma}\label{lem:rank}
If $(C,r)\in W_t$, the cops can force capture in at most $t$ cop moves.
If $L$ is a set of legal states such that every cop configuration has a
state in $L$, and for each $(C,r)\in L$ and each $D$ with
$C\leadsto D$ there is an $r'\in R_D(r)$ for which $(D,r')\in L$, then
the robber can evade $k$ cops forever.
\end{lemma}

\begin{proof}
The first assertion is induction on the least $t$ for which the state enters
the attractor.  A witnessing move either captures immediately or leaves only
responses of smaller rank.  For the second assertion, the robber initially
chooses a certified state in $L$, then uses the closure condition after each
cop move.  Induction over the rounds keeps the state in $L$, so capture never
occurs.
\end{proof}

The certificate represents a cop configuration by its lexicographically
ordered multiset.  It does not quotient by any graph automorphism; the only
quotient is permutation of cop labels.  Coordinates are encoded by
$(x,y)\mapsto4x+y$, for $x,y\in\{0,1,2,3\}$.  The 32 undirected edges are
stored as sorted integer pairs.  For an unambiguous graph check, the compact
JSON serialization
\[
\texttt{\{"vertices":[0,...,15],"edges":[...sorted pairs...]\}}
\]
has SHA-256 digest
\begin{center}
\small\nolinkurl{dc5ff1ba83c901714eae763d51d3c87d763524cff2a1b8c4a57929254d106d5e}
\end{center}
The accompanying compact JSON attractor certificate has SHA-256 digest
\begin{center}
\small\nolinkurl{a544375352b451c4f703477b9ea88c3dfad6096573cb3c1b4e2cc80dbb8bc779}
\end{center}

For four cops there are 3,876 multiset configurations and 48,960 legal states.
Every legal state belongs to $W_4$.  For each state, the certificate records
its first-entry rank and a witnessing destination $D$.  The rank histogram is
\[
\begin{array}{c|rrrr}
\text{rank}&1&2&3&4\\ \hline
\text{states}&32{,}944&9{,}344&6{,}496&176.
\end{array}
\]
The verifier checks that each stored destination lifts to a legal move of the
labelled cops.  It then checks either immediate capture or, for every vertex
in $R_D(r)$, a positive rank strictly smaller than the source rank.  Thus
Lemma~\ref{lem:rank} supplies a replayable four-cop strategy from every
placement and guarantees capture within four cop moves.

For three cops there are 816 multiset configurations and 10,880 legal states.
The fixed attractor contains 6,304 states, and its complement $L$ contains
4,576.  The certificate verifies that $L$ is nonempty above every one of the
816 initial configurations.  For each $(C,r)\in L$ and every distinct
destination multiset $D$, it stores a vertex $r'\in R_D(r)$ with
$(D,r')\in L$.  In particular, no such cop move captures $r$.  The second
part of Lemma~\ref{lem:rank} therefore gives a robber strategy against three
cops.

The separately implemented checker constructs adjacency directly from
modular coordinate differences and imports no search or exporter code.  It
expands the certificate to all 810,000 legal ordered four-cop states.  For the
lower bound it checks 2,662,000 ordered cop moves from the 21,296 labelled
states represented by $L$, including all $5^3$ stays-or-steps from each
labelled three-cop tuple.  All comparisons use integers, finite sets, and bit
masks.  Combining the verified upper and lower bounds proves
Theorem~\ref{thm:main}.

\section{Reproducibility and disclosure}

The ancillary files include the certificate exporter, the compact certificate,
the separately written labelled-cop checker, and its deterministic report.
From the project root, the decisive checks are reproduced by
\begin{verbatim}
python3 evidence/export_attractor_certificate.py
python3 evidence/check_labelled_certificate.py
\end{verbatim}
The first command reconstructs the certificate from the game definition; the
second checks the certificate after reconstructing the graph and all labelled
moves independently.  Neither command uses floating-point arithmetic, random
seeds, numerical tolerances, or an external solver.

OpenAI Codex assisted with literature searching, finite-state design, draft
code, adversarial audit organization, and initial manuscript preparation.
Language-model output is not used as mathematical evidence.  The theorem is
supported by the explicit attractor argument, the exact certificate, and the
deterministic verifier described above.  No proof assistant was used.

\bibliographystyle{plain}
\bibliography{references}

\end{document}
